% ============================================================
% Question Bank: ch09-04 Probability Models
% Topic Title: Probability Models
% CCSS: 7.SP.C.7
% Questions: 30 (18 MC + 7 SA + 5 Graphical)
% ============================================================

% --- Multiple Choice Questions (18) ---

\begin{practiceQuestion}{9-4-q01}{mc}
\begin{questionText}
Which of the following is an example of a uniform probability model?
\end{questionText}
\choiceA{A bag with $3$ red, $5$ blue, and $2$ green marbles}
\choiceB{A fair number cube with $6$ sides}
\choiceC{A spinner with sections of different sizes}
\choiceD{A weighted coin}
\correctAnswer{B}
\explanation{A fair number cube has $6$ equally likely outcomes, making it a uniform model.}
\end{practiceQuestion}

\begin{practiceQuestion}{9-4-q02}{mc}
\begin{questionText}
In a uniform probability model, what is true about the probability of each outcome?
\end{questionText}
\choiceA{Some outcomes are more likely than others}
\choiceB{All outcomes have the same probability}
\choiceC{One outcome always has probability $1$}
\choiceD{The probabilities do not add up to $1$}
\correctAnswer{B}
\explanation{In a uniform model, all outcomes are equally likely.}
\end{practiceQuestion}

\begin{practiceQuestion}{9-4-q03}{mc}
\begin{questionText}
A bag contains $4$ red marbles and $1$ blue marble. Which type of probability model does this represent?
\end{questionText}
\choiceA{Uniform}
\choiceB{Non-uniform}
\choiceC{Both uniform and non-uniform}
\choiceD{Neither}
\correctAnswer{B}
\explanation{The probability of drawing red ($\frac{4}{5}$) is different from drawing blue ($\frac{1}{5}$), so this is a non-uniform model.}
\end{practiceQuestion}

\begin{practiceQuestion}{9-4-q04}{mc}
\begin{questionText}
A spinner has $3$ equal sections: red, blue, and green. What is the probability of landing on any one color?
\end{questionText}
\choiceA{$\frac{1}{2}$}
\choiceB{$\frac{1}{4}$}
\choiceC{$\frac{1}{3}$}
\choiceD{$\frac{2}{3}$}
\correctAnswer{C}
\explanation{With $3$ equal sections, each has probability $\frac{1}{3}$. This is a uniform model.}
\end{practiceQuestion}

\begin{practiceQuestion}{9-4-q05}{mc}
\begin{questionText}
In any probability model, what must the sum of all outcome probabilities equal?
\end{questionText}
\choiceA{$0$}
\choiceB{$0.5$}
\choiceC{$1$}
\choiceD{It depends on the model}
\correctAnswer{C}
\explanation{In any probability model (uniform or non-uniform), all probabilities must add up to $1$.}
\end{practiceQuestion}

\begin{practiceQuestion}{9-4-q06}{mc}
\begin{questionText}
A spinner is divided into $4$ sections with probabilities: red $= 0.4$, blue $= 0.3$, green $= 0.2$, yellow $= 0.1$. What type of model is this?
\end{questionText}
\choiceA{Uniform, because there are $4$ sections}
\choiceB{Non-uniform, because the probabilities are different}
\choiceC{Invalid, because the probabilities don't add up to $1$}
\choiceD{Uniform, because all sections are on the same spinner}
\correctAnswer{B}
\explanation{The probabilities are $0.4, 0.3, 0.2, 0.1$ — all different, so this is non-uniform. They sum to $1.0$, making it valid.}
\end{practiceQuestion}

\begin{practiceQuestion}{9-4-q07}{mc}
\begin{questionText}
A coin is weighted so that heads comes up $60\%$ of the time. What is the probability of tails?
\end{questionText}
\choiceA{$0.60$}
\choiceB{$0.50$}
\choiceC{$0.40$}
\choiceD{$0.30$}
\correctAnswer{C}
\explanation{$P(\text{tails}) = 1 - 0.60 = 0.40$. Since heads and tails have different probabilities, this is a non-uniform model.}
\end{practiceQuestion}

\begin{practiceQuestion}{9-4-q08}{mc}
\begin{questionText}
Which of the following sets of probabilities could represent a valid probability model with $3$ outcomes?
\end{questionText}
\choiceA{$0.2, 0.5, 0.4$}
\choiceB{$0.3, 0.3, 0.4$}
\choiceC{$0.1, 0.2, 0.5$}
\choiceD{$0.5, 0.5, 0.5$}
\correctAnswer{B}
\explanation{$0.3 + 0.3 + 0.4 = 1.0$. The other options do not sum to $1$.}
\end{practiceQuestion}

\begin{practiceQuestion}{9-4-q09}{mc}
\begin{questionText}
A teacher predicted that a spinner would land on red $\frac{1}{4}$ of the time. After $80$ spins, red appeared $25$ times. How do the observed and predicted results compare?
\end{questionText}
\choiceA{They are very different — the model is wrong}
\choiceB{The observed ($\frac{25}{80} = 0.3125$) is close to the predicted ($0.25$)}
\choiceC{The observed result is exactly equal to the predicted result}
\choiceD{The predicted probability is higher than the observed}
\correctAnswer{B}
\explanation{Predicted: $\frac{1}{4} = 0.25$. Observed: $\frac{25}{80} = 0.3125$. They are close, which supports the model.}
\end{practiceQuestion}

\begin{practiceQuestion}{9-4-q10}{mc}
\begin{questionText}
A fair die is rolled. Which probability model correctly shows the probability of each outcome?
\end{questionText}
\choiceA{$P(1) = \frac{1}{6}$, $P(2) = \frac{1}{6}$, $P(3) = \frac{1}{6}$, $P(4) = \frac{1}{6}$, $P(5) = \frac{1}{6}$, $P(6) = \frac{1}{6}$}
\choiceB{$P(1) = \frac{1}{6}$, $P(2) = \frac{2}{6}$, $P(3) = \frac{3}{6}$, $P(4) = 0$, $P(5) = 0$, $P(6) = 0$}
\choiceC{$P(\text{even}) = \frac{1}{2}$, $P(\text{odd}) = \frac{1}{2}$}
\choiceD{$P(1) = \frac{1}{3}$, $P(2) = \frac{1}{3}$, $P(3) = \frac{1}{3}$}
\correctAnswer{A}
\explanation{A fair die has $6$ equally likely outcomes, each with probability $\frac{1}{6}$.}
\end{practiceQuestion}

\begin{practiceQuestion}{9-4-q11}{mc}
\begin{questionText}
A bag contains $5$ yellow and $5$ purple marbles. Lisa picks a marble, records the color, and replaces it $40$ times. She gets yellow $24$ times. Which statement is correct?
\end{questionText}
\choiceA{The model predicts $P(\text{yellow}) = 0.5$, and the experiment gives $P(\text{yellow}) = 0.6$, so there is a small difference}
\choiceB{The model is wrong because $24 \neq 20$}
\choiceC{The model predicts $P(\text{yellow}) = 0.6$}
\choiceD{The experiment proves yellow is more likely than purple}
\correctAnswer{A}
\explanation{Predicted: $\frac{5}{10} = 0.5$. Experimental: $\frac{24}{40} = 0.6$. The small difference is expected from natural variation.}
\end{practiceQuestion}

\begin{practiceQuestion}{9-4-q12}{mc}
\begin{questionText}
Why might you use a non-uniform model instead of a uniform model?
\end{questionText}
\choiceA{Because non-uniform models are always more accurate}
\choiceB{Because the outcomes are not all equally likely}
\choiceC{Because uniform models cannot be used in real life}
\choiceD{Because non-uniform models always sum to more than $1$}
\correctAnswer{B}
\explanation{Non-uniform models are used when outcomes have different probabilities. For example, a weighted spinner or a biased coin.}
\end{practiceQuestion}

\begin{practiceQuestion}{9-4-q13}{mc}
\begin{questionText}
A probability model has outcomes A, B, and C with $P(A) = 0.5$ and $P(B) = 0.3$. What is $P(C)$?
\end{questionText}
\choiceA{$0.8$}
\choiceB{$0.3$}
\choiceC{$0.2$}
\choiceD{$0.1$}
\correctAnswer{C}
\explanation{All probabilities must sum to $1$. $P(C) = 1 - 0.5 - 0.3 = 0.2$.}
\end{practiceQuestion}

\begin{practiceQuestion}{9-4-q14}{mc}
\begin{questionText}
A student creates a probability model for a spinner: $P(\text{red}) = 0.25$, $P(\text{blue}) = 0.25$, $P(\text{green}) = 0.25$, $P(\text{yellow}) = 0.25$. After $200$ spins, the results are: red $= 55$, blue $= 48$, green $= 52$, yellow $= 45$. Is the model a good fit?
\end{questionText}
\choiceA{No, the model is completely wrong}
\choiceB{Yes, the observed frequencies are reasonably close to $50$ each}
\choiceC{No, because the results are not exactly $50$ each}
\choiceD{Yes, because $200$ is divisible by $4$}
\correctAnswer{B}
\explanation{Each color should appear about $50$ times ($200 \times 0.25$). The results ($55, 48, 52, 45$) are all close to $50$, supporting the model.}
\end{practiceQuestion}

\begin{practiceQuestion}{9-4-q15}{mc}
\begin{questionText}
Which of these is NOT a valid probability model?
\end{questionText}
\choiceA{$P(A) = 0.5$, $P(B) = 0.5$}
\choiceB{$P(A) = 0.7$, $P(B) = 0.2$, $P(C) = 0.1$}
\choiceC{$P(A) = 0.6$, $P(B) = 0.6$}
\choiceD{$P(A) = 1$}
\correctAnswer{C}
\explanation{$0.6 + 0.6 = 1.2 > 1$. The probabilities must sum to exactly $1$.}
\end{practiceQuestion}

\begin{practiceQuestion}{9-4-q16}{mc}
\begin{questionText}
A game has $4$ outcomes: Win Big, Win Small, Lose Small, Lose Big. The probabilities are $0.05$, $0.25$, $0.45$, and $0.25$. Which outcome is LEAST likely?
\end{questionText}
\choiceA{Win Big}
\choiceB{Win Small}
\choiceC{Lose Small}
\choiceD{Lose Big}
\correctAnswer{A}
\explanation{Win Big has the smallest probability ($0.05$), making it least likely.}
\end{practiceQuestion}

\begin{practiceQuestion}{9-4-q17}{mc}
\begin{questionText}
A bag has tiles labeled A, A, A, B, B. What is the probability of drawing a B?
\end{questionText}
\choiceA{$\frac{1}{5}$}
\choiceB{$\frac{2}{3}$}
\choiceC{$\frac{2}{5}$}
\choiceD{$\frac{3}{5}$}
\correctAnswer{C}
\explanation{There are $2$ B tiles out of $5$ total. $P(B) = \frac{2}{5}$.}
\end{practiceQuestion}

\begin{practiceQuestion}{9-4-q18}{mc}
\begin{questionText}
Ben predicts that a coin will land on heads $50\%$ of the time. He flips the coin $20$ times and gets $14$ heads. What should he conclude?
\end{questionText}
\choiceA{The coin is definitely biased}
\choiceB{His model is exactly correct}
\choiceC{The result is somewhat different from the prediction, but $20$ trials is a small sample}
\choiceD{He should change his model to predict $70\%$ heads}
\correctAnswer{C}
\explanation{$\frac{14}{20} = 0.70$ vs. predicted $0.50$. With only $20$ trials, variation is expected. More trials would provide a better comparison.}
\end{practiceQuestion}

% --- Short Answer Questions (7) ---

\begin{practiceQuestion}{9-4-q19}{sa}
\begin{questionText}
A probability model has $5$ equally likely outcomes. What is the probability of each outcome? Write your answer as a fraction.
\end{questionText}
\correctAnswer{$\frac{1}{5}$}
\explanation{In a uniform model with $5$ outcomes, each has probability $\frac{1}{5}$.}
\end{practiceQuestion}

\begin{practiceQuestion}{9-4-q20}{sa}
\begin{questionText}
A probability model has outcomes X, Y, and Z. $P(X) = 0.15$ and $P(Y) = 0.45$. What is $P(Z)$?
\end{questionText}
\correctAnswer{$0.4$}
\explanation{$P(Z) = 1 - 0.15 - 0.45 = 0.40$.}
\end{practiceQuestion}

\begin{practiceQuestion}{9-4-q21}{sa}
\begin{questionText}
A bag has $6$ red, $4$ blue, and $10$ white marbles. Write the probability model by listing $P(\text{red})$, $P(\text{blue})$, and $P(\text{white})$ as fractions in simplest form.
\end{questionText}
\correctAnswer{$P(\text{red}) = \frac{3}{10}$, $P(\text{blue}) = \frac{1}{5}$, $P(\text{white}) = \frac{1}{2}$}
\explanation{Total: $20$. $P(\text{red}) = \frac{6}{20} = \frac{3}{10}$, $P(\text{blue}) = \frac{4}{20} = \frac{1}{5}$, $P(\text{white}) = \frac{10}{20} = \frac{1}{2}$.}
\end{practiceQuestion}

\begin{practiceQuestion}{9-4-q22}{sa}
\begin{questionText}
Is flipping a fair coin a uniform or non-uniform probability model?
\end{questionText}
\correctAnswer{Uniform}
\explanation{A fair coin has $2$ equally likely outcomes (heads, tails), each with probability $\frac{1}{2}$. This is uniform.}
\end{practiceQuestion}

\begin{practiceQuestion}{9-4-q23}{sa}
\begin{questionText}
A spinner model predicts $P(\text{blue}) = 0.30$. In $50$ spins, blue appears $18$ times. What is the experimental probability? Is it close to the predicted value?
\end{questionText}
\correctAnswer{$0.36$; yes, it is close}
\explanation{$P = \frac{18}{50} = 0.36$. The predicted value is $0.30$. The difference ($0.06$) is small and expected from natural variation.}
\end{practiceQuestion}

\begin{practiceQuestion}{9-4-q24}{sa}
\begin{questionText}
Create a non-uniform probability model for a spinner with $3$ sections (A, B, C) where A is twice as likely as B, and C has the same probability as B. Write the probability of each section.
\end{questionText}
\correctAnswer{$P(A) = \frac{1}{2}$, $P(B) = \frac{1}{4}$, $P(C) = \frac{1}{4}$}
\explanation{Let $P(B) = x$. Then $P(A) = 2x$ and $P(C) = x$. Total: $2x + x + x = 4x = 1$, so $x = \frac{1}{4}$. $P(A) = \frac{1}{2}$, $P(B) = P(C) = \frac{1}{4}$.}
\end{practiceQuestion}

\begin{practiceQuestion}{9-4-q25}{sa}
\begin{questionText}
A probability model assigns $P(A) = 0.35$, $P(B) = 0.25$, $P(C) = 0.25$, and $P(D) = 0.20$. Is this a valid model? Explain.
\end{questionText}
\correctAnswer{No}
\explanation{$0.35 + 0.25 + 0.25 + 0.20 = 1.05$, which is greater than $1$. A valid model must have probabilities that sum to exactly $1$.}
\end{practiceQuestion}

% --- Graphical Questions (3 GMC + 2 GSA) ---

\begin{practiceQuestion}{9-4-q26}{gmc}
\begin{questionText}
Look at the spinner below. Which probability model correctly represents this spinner?

\begin{center}
\begin{tikzpicture}[scale=1.2]
  \fill[funRed!50] (0,0) -- (0:1.5) arc (0:180:1.5) -- cycle;
  \fill[funBlue!50] (0,0) -- (180:1.5) arc (180:270:1.5) -- cycle;
  \fill[funGreen!50] (0,0) -- (270:1.5) arc (270:360:1.5) -- cycle;
  \draw[thick] (0,0) circle (1.5);
  \draw[thick] (0,0) -- (0:1.5);
  \draw[thick] (0,0) -- (180:1.5);
  \draw[thick] (0,0) -- (270:1.5);
  \node at (90:0.8) {\textbf{Red}};
  \node at (225:0.8) {\textbf{Blue}};
  \node at (315:0.8) {\textbf{Green}};
\end{tikzpicture}
\end{center}
\end{questionText}
\choiceA{$P(\text{Red}) = \frac{1}{3}$, $P(\text{Blue}) = \frac{1}{3}$, $P(\text{Green}) = \frac{1}{3}$}
\choiceB{$P(\text{Red}) = \frac{1}{2}$, $P(\text{Blue}) = \frac{1}{4}$, $P(\text{Green}) = \frac{1}{4}$}
\choiceC{$P(\text{Red}) = \frac{1}{4}$, $P(\text{Blue}) = \frac{1}{4}$, $P(\text{Green}) = \frac{1}{2}$}
\choiceD{$P(\text{Red}) = \frac{1}{2}$, $P(\text{Blue}) = \frac{1}{2}$, $P(\text{Green}) = 0$}
\correctAnswer{B}
\explanation{Red covers half the spinner ($180\degree$), while Blue and Green each cover a quarter ($90\degree$ each). This is a non-uniform model.}
\end{practiceQuestion}

\begin{practiceQuestion}{9-4-q27}{gmc}
\begin{questionText}
The table shows a probability model for a spinner. The spinner is spun $100$ times and the results are shown. Does the observed data support the model?

\begin{center}
\begin{tabular}{|c|c|c|}
\hline
\textbf{Color} & \textbf{Model $P$} & \textbf{Observed Frequency} \\
\hline
Red & $0.40$ & $38$ \\
\hline
Blue & $0.35$ & $36$ \\
\hline
Yellow & $0.25$ & $26$ \\
\hline
\end{tabular}
\end{center}
\end{questionText}
\choiceA{No, the data is very different from the model}
\choiceB{Yes, the observed frequencies are close to the predicted values}
\choiceC{No, because Red appeared fewer times than predicted}
\choiceD{The model cannot be checked with only $100$ spins}
\correctAnswer{B}
\explanation{Predicted: Red $= 40$, Blue $= 35$, Yellow $= 25$. Observed: $38$, $36$, $26$. These are close, supporting the model.}
\end{practiceQuestion}

\begin{practiceQuestion}{9-4-q28}{gmc}
\begin{questionText}
Look at the two bags below. Which bag represents a uniform probability model for drawing a marble?

\begin{center}
\begin{tikzpicture}[scale=0.65]
  % Bag 1
  \draw[thick,rounded corners=8pt] (-3,-1.5) rectangle (0,1.5);
  \node[above,font=\bfseries] at (-1.5,1.5) {Bag 1};
  \fill[funRed] (-2.3,0.5) circle (0.3);
  \fill[funRed] (-1.5,0.5) circle (0.3);
  \fill[funBlue] (-0.7,0.5) circle (0.3);
  \fill[funBlue] (-2.3,-0.5) circle (0.3);
  \fill[funGreen] (-1.5,-0.5) circle (0.3);
  \fill[funGreen] (-0.7,-0.5) circle (0.3);
  % Bag 2
  \draw[thick,rounded corners=8pt] (1,-1.5) rectangle (4,1.5);
  \node[above,font=\bfseries] at (2.5,1.5) {Bag 2};
  \fill[funRed] (1.7,0.5) circle (0.3);
  \fill[funRed] (2.5,0.5) circle (0.3);
  \fill[funRed] (3.3,0.5) circle (0.3);
  \fill[funRed] (1.7,-0.5) circle (0.3);
  \fill[funBlue] (2.5,-0.5) circle (0.3);
  \fill[funGreen] (3.3,-0.5) circle (0.3);
\end{tikzpicture}
\end{center}
\end{questionText}
\choiceA{Bag 1 only}
\choiceB{Bag 2 only}
\choiceC{Both bags}
\choiceD{Neither bag}
\correctAnswer{A}
\explanation{Bag 1 has $2$ of each color ($2$ red, $2$ blue, $2$ green), so each color is equally likely — uniform. Bag 2 has $4$ red, $1$ blue, $1$ green — non-uniform.}
\end{practiceQuestion}

\begin{practiceQuestion}{9-4-q29}{gsa}
\begin{questionText}
The bar graph shows the predicted number of times each color should appear if a spinner is spun $60$ times. Based on this model, what is the probability of landing on green? Write your answer as a fraction in simplest form.

\begin{center}
\begin{tikzpicture}[scale=0.55]
  \draw[thick,->] (0,0) -- (0,32) node[above,font=\small] {Predicted};
  \draw[thick,->] (0,0) -- (7,0);
  \foreach \y in {5,10,...,30} {
    \draw[gray!30] (0,\y) -- (6.5,\y);
    \node[left,font=\small] at (0,\y) {\y};
  }
  \fill[funRed!70] (0.5,0) rectangle (1.5,30);
  \fill[funBlue!70] (2.5,0) rectangle (3.5,18);
  \fill[funGreen!70] (4.5,0) rectangle (5.5,12);
  \node[below,font=\small] at (1,0) {Red};
  \node[below,font=\small] at (3,0) {Blue};
  \node[below,font=\small] at (5,0) {Green};
\end{tikzpicture}
\end{center}
\end{questionText}
\correctAnswer{$\frac{1}{5}$}
\explanation{Green is predicted to appear $12$ out of $60$ times. $P = \frac{12}{60} = \frac{1}{5}$.}
\end{practiceQuestion}

\begin{practiceQuestion}{9-4-q30}{gsa}
\begin{questionText}
The table below shows a probability model for a game spinner. One probability is missing. Find the missing probability for outcome D.

\begin{center}
\begin{tabular}{|c|c|}
\hline
\textbf{Outcome} & \textbf{Probability} \\
\hline
A & $0.30$ \\
\hline
B & $0.25$ \\
\hline
C & $0.15$ \\
\hline
D & ? \\
\hline
\end{tabular}
\end{center}
\end{questionText}
\correctAnswer{$0.30$}
\explanation{$P(D) = 1 - (0.30 + 0.25 + 0.15) = 1 - 0.70 = 0.30$.}
\end{practiceQuestion}
